% =============================================================
% frontmatter/introduction.tex
% Dissertation-level Introduction: Long-Form Synthesis Abstract
% Evidence base: Abstracts, Introductions, and Conclusions of
% Essays 1, 2, and 3.
% =============================================================

\chapter*{Introduction}
\addcontentsline{toc}{chapter}{Introduction}
\markboth{Introduction}{Introduction}

\section*{The Problem of Productive Capacity in Political Economy}

Across both orthodox and heterodox traditions in macroeconomics, the rate of capacity utilization ($\mu \equiv Y / Y^p$) occupies an indispensable theoretical position. Central banks monitor it as an inflation sentinel and a gauge of cyclical slack. Post-Keynesian and Kaleckian growth models treat it as the central adjusting variable that absorbs discrepancies between autonomous demand and productive capacity. Sraffian and classical-Marxian frameworks examine it to understand the gravitation of market prices toward prices of production, the distribution of surplus value, and the secular trajectory of the general rate of profit. Yet despite its pervasive appearance in empirical equations and policy discourse, the denominator of this ratio---productive capacity ($Y^p$)---is an unobserved quantity. Statistical agencies measure realized physical output, and national income accountants record capital expenditures and depreciated assets. No agency records the potential capacity of the capital stock directly.

Every claim about capacity utilization is therefore fundamentally a claim about an unobserved construct, resting entirely on the conventions deployed to locate its denominator. In empirical macroeconomics, this denominator has historically been supplied through three conventional strategies, all of which conceal the underlying transformation of accumulation into potential output. The first strategy treats utilization as an observed macroeconomic primitive, relying on managerial survey responses such as those compiled by the Federal Reserve Board or the McGraw-Hill plant-and-equipment surveys. Yet managerial surveys reflect subjective operating preferences and administrative forecasting routines rather than objective physical or economic frontiers. The second strategy anchors capacity to an inherited engineering maximum or a fixed, convention-bound normal operating rate imported from outside the national accounting framework. The third and most pervasive approach recovers capacity as an empirical residual, fitting a statistical trend or a smooth aggregate production function through observed output conditional on capital stock. 

This third approach introduces a structural identity trap. By fitting output to capital under the presumption of smooth, balanced growth, conventional filtering and cointegration exercises constrain the long-run capacity-to-capital elasticity ($\theta \equiv \partial \ln Y^p / \partial \ln K$) to unity by construction. Potential capacity is assumed to expand in permanent lockstep with accumulated capital ($\theta = 1$). Utilization is then measured merely as the transient fluctuation of realized demand around a benchmark that assumed away the very economic relation required to identify it. When utilization is recovered by assumption, it ceases to describe the distance between actual production and potential capacity; it measures only the residual of an unexamined premise.


\section*{The Unifying Analytical Thread: From Accumulation to Capacity Formation}

The unifying analytical problem that binds the three essays of this dissertation together is the theoretical and empirical reconstruction of capacity formation as a historically contingent, institutionally contested, and distributionally mediated process. Rather than treating productive capacity as a friction-free technological ceiling or an engineering given, this dissertation conceptualizes capacity utilization as something an economy \emph{produces} through decentralized investment decisions, labor-process conflicts, and open-economy balance-of-payments constraints.

To move from accumulated capital ($K$) to usable productive capacity ($Y^p$) requires a determinate transformation relation. In a decentralized capitalist economy, private investment commitments do not coordinate into proportional, balanced expansion; they generate disproportionalities and capacity mismatches as a structural norm rather than an accidental friction. Furthermore, the capitalist enterprise responsible for executing this conversion is internally divided. The organization that transforms financial investment into physical machinery and operating capacity cannot maximize managerial control and labor cooperation simultaneously. Instead, it settles into organizational arrangements that sustain production without extinguishing the underlying antagonism between labor and capital. 

Consequently, the transformation elasticity $\theta$ is an empirical object with a historically determinate value that must be estimated rather than posited. Across the three essays, this inquiry unfolds through a deliberate progression across levels of abstraction, historical periods, and geographic positions within the world economy:
\begin{enumerate}[label=(\roman*)]
  \item In the advanced capitalist core, the conversion of accumulation into capacity is mediated by the intensity of exploitation within the labor process and the institutional distribution of income between wages, profits, and intangible capital.
  \item When the analysis is extended to the dependent periphery, that same conversion encounters an external ceiling: domestic capital accumulation is structurally import-dependent, meaning that the incentive to mechanize is severed from the physical means to import capital goods whenever foreign-exchange reserves run out.
  \item At the highest frequency of crisis, this external boundary materializes on the balance sheet of the peripheral central bank, where the structural subordination of dependent economies within an asymmetric international currency hierarchy compels the monetary authority to monetize enterprise deficits defensively, resolving acute balance-of-payments shocks through endogenous accommodation rather than autonomous monetary policy.
\end{enumerate}


\section*{Overarching Contributions of the Dissertation}

Synthesizing the empirical and theoretical findings of the three essays, the overarching contributions of the dissertation span four distinct dimensions:

\subsection*{Theoretical Contribution}
The dissertation reframes capacity utilization from a reported macroeconomic primitive into a \emph{derived residual} of an underlying transformation process. Joining classical-Marxian theories of choice of technique and decentralized accumulation with the theory of the internally divided firm and structuralist dependency theory, the dissertation demonstrates that the long-run elasticity linking capital accumulation to potential capacity formation ($\theta$) systematically deviates from the balanced-growth benchmark ($\theta = 1$). In the core, this deviation manifests as a structural overaccumulation regime ($\theta < 1$) where capital accumulation outpaces capacity formation. In the periphery, the transformation relation is non-linear and state-dependent, exhibiting an endogenous regime switch governed by external solvency constraints.

\subsection*{Measurement and Methodological Contribution}
The dissertation contributes a rigorous econometric and forensic measurement framework to applied political economy. It reconstructs long-span corporate and national capital stocks using the Generalized Perpetual Inventory Method (GPIM) with finite service lives and Weibull physical retirement functions, correcting for imputed interest and supervisory compensation distortions. It establishes the econometric fragility of single-equation cointegration models through combinatorial specification grids, demonstrates the necessity of system-level cointegration (VECM) conditioned on distributive conflict, implements state-dependent cointegrating regressions with endogenously dated structural breaks (IM-OLS with fixed-$b$ inference), and formalizes the Central Bank Solvency Ratio and Solvency Growth Gap within non-linear Threshold Vector Autoregressions (TVAR) with Generalized Impulse Response Functions (GIRFs).

\subsection*{Empirical Contribution}
The empirical investigations yield robust, calibrated parameters across the United States and Chile. For the post-war US corporate sector (1947--2011 and 1931--2024), the transformation elasticity is strictly below unity across all specifications, drifting moderately from $0.37$ in 1960 to $0.43$ in 2024 as investment migrated into intellectual property and effective wage pressure declined. For Chile (1943--2010), a composite balance-of-payments threshold at $\hat{\gamma} = 0.425$ switches off the mechanization response: the wage-share elasticity of the capital composition gap collapses from $+3.29$ in the unconstrained state to $+0.14$ under foreign-exchange rationing. At monthly frequency during the Unidad Popular (1970--1973), crossing the central-bank insolvency threshold ($\Delta s_{t-1} \le -4.615\%$) establishes that reverse monetary accommodation from inflation to money creation ($\pi \to g_H$, 10-month persistence, 6.54 percentage points cumulative) dominates forward transmission ($g_H \to \pi$, 9-month persistence, 3.55 percentage points cumulative).

\subsection*{Comparative and Historical Contribution}
The dissertation builds an integrated historical bridge connecting the long-run trajectory of accumulation in the advanced capitalist center with peripheral vulnerability in Latin America. It historicizes the shift of capacity metrics from post-war demand-management tools to neoliberal output-gap inflation sentinels, analyzes how intangible asset accumulation reconstituted US corporate profitability after Fordism, and provides an archival, balance-sheet reconstruction of Chile's socialist transition that refutes monocausal monetarist narratives of unforced populist mismanagement.


\section*{Essay 1: Critical Replication of Shaikh's Capacity Utilization Measure}

The first essay investigates whether the long-run output--capital relation can serve as an autonomous, accounting-based foundation for measuring productive capacity, or whether treating this relation as an engineering primitive masks the historical conflicts that govern accumulation. \citet{Shaikh2016} proposed a portable, single-equation autoregressive distributed lag (ARDL) cointegration strategy to locate the productive capacity ceiling for the US corporate sector (1947--2011), arguing that the long-run output-to-capital coefficient reflects a stable technological relation. 

Essay~1 subjects this strategy to a forensic critical replication and econometric stress test. Re-examining the single-equation ARDL specification across a combinatorial grid of 500 admissible models reveals what the essay terms ``disciplined non-uniqueness'': rather than converging on a unique technical parameter, the recovered transformation elasticity ranges from 0.65 to 0.95, varying systematically with arbitrary lag selections and historical dummy variables. 

To evaluate whether this relation holds as a genuine long-run attractor, Essay~1 shifts the analysis to a system-level Vector Error Correction Model (VECM). In a bivariate system consisting solely of output and capital stock, the relation fails cointegration and dynamic stability tests across all lag profiles and deterministic configurations. Cointegration emerges exclusively in a trivariate VECM when the rate of exploitation ($e_t \equiv \pi_t / \omega_t$, the ratio of the profit share to the wage share) enters the state vector. Once trend-containing specifications are excluded on grounds of economic admissibility, the surviving cointegrating vector yields an estimated transformation elasticity strictly below unity ($\hat{\theta} < 1.0$), with a point estimate of $\theta \approx 0.82$. 

This finding establishes that the US corporate sector operated in a structural overaccumulation regime where capital growth outstripped the expansion of productive capacity. Moreover, system speed-of-adjustment parameters identify capital stock as weakly exogenous ($\alpha_k \approx 0$), indicating that functional distribution acts as the primary error-correction channel restoring long-run macroeconomic reproduction. Essay~1 establishes the foundational thesis of the dissertation: productive capacity cannot be estimated in isolation from the distributive conflicts of the labor process, and the transformation elasticity linking capital to capacity must be treated as structurally open and historically conditioned.


\section*{Essay 2: The Transformation of Accumulation into Productive Capacities}

The second essay takes the core empirical finding of Essay~1---that the output--capital relation is distributionally conditioned and non-unitary---and provides it with an analytical foundation and a comparative empirical evaluation across the global divide between center and periphery. If capacity formation does not follow capital accumulation one-for-one, what determines the elasticity that governs their conversion?

Essay~2 answers this question by synthesizing the classical choice-of-technique framework \citep{Kurz1986} with \citeauthor{Okishio1961}'s (\citeyear{Okishio1961,Okishio2022}) model of decentralized investment disproportionality and \citeauthor{Vidal2019}'s (\citeyear{Vidal2019,Vidal2022}) organizational theory of the internally divided firm. Under decentralized accumulation, firms choose techniques to minimize production costs under prevailing wage pressure and competitive threats. Because the organization converting investment into capacity cannot resolve workplace conflict, the conversion rate $\theta \equiv \partial \ln Y^p / \partial \ln K$ is structurally state-dependent.

The essay estimates this transformation elasticity across two long historical panels. For the United States (1931--2024), it deploys the Integrated Modified OLS (IM-OLS) estimator of \citet{VogelsangWagner2024} with fixed-$b$ inference and endogenously dated structural breaks \citep{BaiPerron2003}. It decomposes the capital stock into machinery, equipment, non-residential structures, and intellectual property products, scaling distributive conflict by effective wage pressure ($C_t = \omega_{t-1}(1 - \iota_{t-1})$). The estimates confirm that $\theta$ remained strictly below unity throughout nine decades, drifting from $0.37$ in 1960 to $0.43$ in 2024. This modest upward drift did not reflect technical efficiency; rather, it was driven by the secular decline of the wage share and the diversion of corporate capital into intellectual property assets, while capacity utilization fell to $93.8\%$ in 2024 from its wartime peak.

To evaluate how this relation operates in the Global South, Essay~2 extends the model to Chile (1943--2010). In the periphery, capital accumulation is bound by the Balance of Payments (BoP) constraint: core machinery and equipment cannot be produced domestically and must be imported using foreign exchange. Using a threshold cointegrating multivariate polynomial regression (TCMPR) anchored in a composite BoP index ($Z_t$) combining reserve cover and the import-financing gap \citep{Hansen2000}, Essay~2 detects an endogenous regime shift at $\hat{\gamma} = 0.425$. In the unconstrained regime, rising wage pressure induces defensive mechanization ($+3.29$). Once the balance-of-payments threshold binds, this mechanization response is throttled ($-3.15$ interaction, net elasticity $+0.14$). Peripheral capitalists retain the economic incentive to substitute capital for labor but lose the physical means to do so because foreign exchange is rationed. Capacity utilization is thus politically mediated in the center and externally truncated in the periphery.


\section*{Essay 3: Peripheral Solvency and Endogenous Monetary Accommodation}

The third essay investigates the high-frequency monetary, financial, and balance-sheet mechanisms through which the external constraint identified in Essay~2 detonates macroeconomic instability during an attempted structural transition. It examines the historical trajectory of Chile's Unidad Popular (UP, 1970--1973) under Salvador Allende, a case frequently cited by orthodox economics as the canonical cautionary tale of ``macroeconomic populism'' and unforced monetary mismanagement \citep{DornbuschEdwards1990,Edwards2024}.

Essay~3 reframes this historical episode by treating dependency theory as a Lakatosian scientific research programme \citep{Lakatos1978,Kvangraven2021}, equipping its protective belt with an open-economy central-bank balance-sheet framework. In an asymmetric international currency hierarchy, national fiat money issued by a peripheral central bank ($H_t$) cannot function as world money ($F_t$). The monetary authority faces what Essay~3 terms the \emph{Peripheral Paradox}: domestic credit creation is naturally endogenous, but international settlement is bound to an exogenously constrained stock of foreign exchange reserves. 

Deploying a multi-frequency empirical design integrating annual series (1920--2010) and continuous monthly archival records from the Central Bank of Chile (1960--1980), the essay tests the causal ordering between monetary emission, inflation, and external reserve depletion. In 1971, Central Bank base money expanded by over 100\% while inflation decelerated from 35\% to 22\%, manufacturing output surged by 12.1\%, and capacity utilization reached an unprecedented peak of 98.4\%. This non-inflationary expansion was made possible because inherited foreign exchange reserves buffered import requirements and mobilized idle capacity. 

The breakdown occurred only after mid-1972, when President Nixon's closure of the gold window destabilized the international monetary system, copper prices declined, and a targeted credit blockade canceled \$220~million in revolving commercial credit lines. To evaluate dynamic shock transmission across solvency states, the essay estimates a non-linear Threshold Vector Autoregression (TVAR) with Generalized Impulse Response Functions (GIRFs), where the transition variable is the predetermined monthly Solvency Growth Gap ($\Delta s_{t-1} \equiv g^F_{t-1} - g^H_{t-1}$). The threshold identifies a structural transition into conditions of central bank insolvency at $\Delta s_{t-1} \le -4.615\%$. Within this insolvency regime, reverse defensive accommodation ($\pi_m \to g^H$) dominates forward monetary transmission ($g^H \to \pi_m$) in both duration (10 consecutive months versus 9) and cumulative magnitude (6.54 percentage points versus 3.55 percentage points). Point-wise multiplier differences confirm that reverse accommodation significantly exceeds forward transmission across nine consecutive months ($h=2,\dots,10$). Rather than acting as an autonomous driver of inflation, the peripheral central bank functioned as an endogenous accommodation valve, forced to monetize enterprise deficits to prevent immediate payroll collapse and distribution breakdown under foreign-exchange starvation.


\section*{Methodological Plurality and the Sequence of the Argument}

The three essays do not deploy identical econometric estimators, nor are they restricted to a single temporal frequency. Methodological plurality is an intentional design feature that serves the common analytical problem across distinct scales of analysis:
\begin{itemize}[leftmargin=1.8em]
  \item In \textbf{Essay~1}, the research problem is forensic and epistemological: evaluating the validity of a widely cited heterodox measurement strategy on its own empirical terrain (the US corporate sector, 1947--2011). This task requires combinatorial single-equation ARDL bounds testing and system-level Johansen VECM cointegration to assess whether an unconditioned output--capital cointegrating vector survives rigorous econometric screening.
  \item In \textbf{Essay~2}, the analytical scale shifts to comparative, long-run structural estimation (US 1931--2024 and Chile 1943--2010). This requires estimators capable of handling non-stationary time series with deterministic shifts, heteroskedasticity, and state-dependent interactions: IM-OLS with fixed-$b$ asymptotic inference and Bai-Perron endogenous break dating for the center, and Hansen profile likelihood threshold cointegration for the periphery.
  \item In \textbf{Essay~3}, the focus shifts from multi-decade secular trends to high-frequency crisis dynamics and directional precedence during a revolutionary rupture (Chile, 1960--1980 at monthly frequency). This requires dynamic time-series methods that can adjudicate temporal ordering and state-dependent feedback without imposing linearity: stationary VAR Granger tests, constraint-based PCMCI causal discovery, and non-linear Threshold Vector Autoregressions with Generalized Impulse Response Functions.
\end{itemize}
This methodological sequence mirrors the conceptual trajectory of the dissertation: from the initial deconstruction of closed-economy capacity measurement (Essay~1), to the structural modeling of state-dependent transformation across center and periphery (Essay~2), to the archival and balance-sheet investigation of high-frequency crisis transmission under external financial subordination (Essay~3).


\section*{Dissertation-Level Implications}

Reading the three essays together reveals a set of fundamental insights into the political economy of capitalist reproduction that remain invisible when the papers are read in isolation.

First, the dissertation dismantles the premise of balanced growth that unites both mainstream output-gap routines and standard heterodox growth models. Whether in the post-war United States or twentieth-century Chile, capital accumulation does not translate into capacity formation one-for-one. The transformation elasticity $\theta$ is strictly non-unitary. In the advanced capitalist core, persistent overaccumulation ($\theta < 1$) imposes a secular downward drag on the capital-capacity ratio ($Y^p/K$). Capital counteracts this drag not through technical stabilization, but by depressing the wage share, intensifying exploitation, and restructuring assets toward intellectual property. 

Second, the dissertation reveals how the international division of labor and the international currency hierarchy fundamentally differentiate the laws of motion of productive capacity across the global divide:
\begin{itemize}[leftmargin=1.8em]
  \item In the \textbf{center}, capacity utilization is mediated primarily by domestic class struggle and realization conditions in the product market. When low unemployment strengthens labor's bargaining position, capital strikes and profitability squeezes induce disinvestment and idle capacity.
  \item In the \textbf{periphery}, capacity utilization is governed from the outside by the balance-of-payments constraint and the availability of world money. Peripheral economies face a double strangulation: they cannot substitute domestic credit for hard currency to purchase imported machinery, and their central bank balance sheets lack the international reserve backing required to insulate domestic price levels from international terms-of-trade shocks.
\end{itemize}

Third, the dissertation reframes the political economy of macroeconomic policy in dependent formations. It establishes that monetary instability in the Global South cannot be diagnosed through the insular categories of closed-economy monetarism or naive chartalist extensions of Modern Monetary Theory. The peripheral state's formal legal authority over domestic fiat currency does not confer command over the universal means of international settlement. When external credit embargoes cut off foreign trade finance, domestic money growth operates as a defensive, accommodating response to enterprise illiquidity rather than an unforced political indulgence. Sustainable progressive transformation in the periphery cannot bypass the structural imperative of external solvency and the material realities of the international currency hierarchy.


\section*{Structure of the Volume}


The remainder of this dissertation is organized into three essays:
\begin{description}[leftmargin=1.8em, font=\normalfont\bfseries]
  \item[Essay 1:] \emph{Critical Replication of Shaikh's Capacity Utilization Measure}.
  \item[Essay 2:] \emph{The Transformation of Accumulation into Productive Capacities: Capacity Utilization in the Center and Periphery}.
  \item[Essay 3:] \emph{Re-visiting the Political Economy of the Rise and Fall of the Unidad Popular: Towards a Global Political Economy Approach}.
\end{description}

\clearpage
{\small
\setlength{\bibsep}{3pt}
\bibliographystyle{apalike}
\bibliography{Chapter1/references,Chapter2/references,Chapter3/references}
}
